% Source: https://artofproblemsolving.com/community/c3h1595782_arithmetic_problem
% Stable identity: aops-topic-1595782-q01
Call a number \emph{good} if it can be written as $\dfrac{m}{3\cdot2^k}$ with $m$ an integer and $k\ge0$ an integer.

\textbf{The starting numbers are good:} $\dfrac{1}{12}=\dfrac{1}{3\cdot2^2}$, $\dfrac{1}{8}=\dfrac{3}{3\cdot2^3}$ and $\dfrac{1}{6}=\dfrac{1}{3\cdot2^1}$.

\textbf{The moves keep numbers good.} Multiplying by $1$ changes nothing, and $\dfrac{1}{2}\cdot\dfrac{m}{3\cdot2^k}=\dfrac{m}{3\cdot2^{k+1}}$. For a sum or difference, write both numbers over the common denominator $3\cdot2^K$, where $K$ is the larger exponent: $\dfrac{m}{3\cdot2^k}\pm\dfrac{m'}{3\cdot2^{k'}}=\dfrac{m2^{K-k}\pm m'2^{K-k'}}{3\cdot2^K}$, which is good.

So every number you can ever obtain is good.

\textbf{$\dfrac{3}{10}$ is not good.} If $\dfrac{3}{10}=\dfrac{m}{3\cdot2^k}$, then $9\cdot2^k=10m$, so $5$ divides $9\cdot2^k$, which is impossible because $9\cdot2^k$ has no prime factor $5$.

Hence $\dfrac{3}{10}$ can never be obtained. $\blacksquare$
